\documentclass[11pt]{amsart}

\usepackage{amsmath,amssymb,amsthm}
\usepackage[margin=1in]{geometry}
\usepackage{mathtools}
\usepackage{stmaryrd}   % \llbracket \rrbracket
\usepackage{booktabs}
\usepackage{hyperref}
\hypersetup{colorlinks=true,linkcolor=blue,citecolor=blue,urlcolor=blue}

%--- theorem environments -------------------------------------------------
\theoremstyle{plain}
\newtheorem{theorem}{Theorem}[section]
\newtheorem{proposition}[theorem]{Proposition}
\newtheorem{lemma}[theorem]{Lemma}
\newtheorem{corollary}[theorem]{Corollary}
\theoremstyle{definition}
\newtheorem{definition}[theorem]{Definition}
\newtheorem{example}[theorem]{Example}
\newtheorem{question}[theorem]{Question}
\theoremstyle{remark}
\newtheorem{remark}[theorem]{Remark}

%--- macros ---------------------------------------------------------------
\newcommand{\C}{\mathcal{C}}
\newcommand{\D}{\mathcal{D}}
\newcommand{\E}{\mathcal{E}}
\newcommand{\Set}{\mathbf{Set}}
\newcommand{\Vect}{\mathbf{Vec}}
\newcommand{\Vecfd}{\mathbf{Vec}_{\mathrm{fd}}}
\newcommand{\Cat}{\mathbf{Cat}}
\newcommand{\Fam}{\mathrm{Fam}}
\newcommand{\Poly}{\mathbf{Poly}}
\newcommand{\Cont}{\mathbf{Cont}}
\newcommand{\Rel}{\mathbf{Rel}}
\newcommand{\Prof}{\mathbf{Prof}}
\newcommand{\Mod}{\mathbf{Mod}}
\newcommand{\Vgen}{\mathcal{V}}
\newcommand{\ext}[1]{\llbracket #1 \rrbracket}
\newcommand{\ihom}[2]{[#1,#2]}
\renewcommand{\lhd}{\vartriangleleft}
\newcommand{\op}{\mathrm{op}}
\newcommand{\id}{\mathrm{id}}
\newcommand{\coprodd}{\textstyle\coprod}
\newcommand{\prodd}{\textstyle\prod}
\DeclareMathOperator{\Nat}{Nat}
\DeclareMathOperator{\Lan}{Lan}
\DeclareMathOperator{\cod}{cod}
\DeclareMathOperator{\dom}{dom}
\DeclareMathOperator{\Sub}{Sub}
\newcommand{\Arr}[1]{#1^{\rightarrow}}
\newcommand{\Exists}{\mathsf{Exists}}
\newcommand{\All}{\mathsf{All}}

\newcommand{\gerund}{}

%==========================================================================
\begin{document}

\title[Approaches to containers over a base]{Approaches to containers over a base:\\
extensivity, cartesian closure, and the linear collapse}

\author{MacBeth}
\address{Kodamai}
\email{neil@kodamai.com}

\date{27 August 2026}

\begin{abstract}
The phrase ``a container in a category'' names at least four inequivalent constructions:
external families over a base $\C$, indexed dependent-polynomial functors over a locally
cartesian closed base, fibrational containers over a comprehension category, and Walker's
locally subcartesian closed weakening. Over $\Set$ they coincide; off $\Set$ they diverge, and
no map has been drawn of \emph{where}. We draw one --- a census of the bases over which the
construction has been run, from toposes down through Weber's and Walker's weakenings to additive
bases --- and argue that two properties of the base --- extensivity
and local cartesian closure --- are the two axes of divergence, and we use the category of vector
spaces, which has neither, as the discriminating example. Only the external-family construction
$\Fam(\C^\op)$ reaches $\Fam(\Vect^\op)$; the other three cannot even form their semantics there.
We make this quantitative with three theorems about $\Fam(\C^\op)$: the container extension is fully
faithful iff the monoidal unit is connected (\emph{not} iff $\C$ is extensive); the Dirichlet tensor
is closed iff a familial-representability condition holds, which over a linear base isolates the
finite-dimensional corner; and the substitution product is left-closed exactly where it collapses
onto the Dirichlet tensor --- so that extensivity, the obstruction to the first two, is what would
have to fail for the third. The upshot is a two-axis classification in which $\Fam(\Vect^\op)$ forces
the external approach. Along the way we develop the fibrational approach beyond its refereeing role
into an explicit \emph{logic of containers}: applying the container construction to the codomain
fibration of $\Set$ yields a bifibration whose fibred quantifiers are the $\mathsf{All}$ and
$\mathsf{Exists}$ predicate liftings, and whose hyperdoctrine is the fibrewise opposite of the usual
one --- container-$\exists$ is $\Set$-$\forall$, and each fibre is a co-topos.
\end{abstract}

\maketitle

%==========================================================================
\section{Introduction}

\subsection{The ambiguity}
A \emph{container} in the sense of Abbott, Altenkirch and Ghani \cite{AAG} is a set $S$ of shapes
together with, for each shape $s$, a set $P_s$ of positions; its extension is the polynomial
endofunctor $\ext{S,P}X = \coprodd_{s\in S} X^{P_s}$ on $\Set$. Containers are the syntax of
strictly positive datatypes, and their theory is by now classical: containers with their morphisms
are equivalent to a full subcategory of $[\Set,\Set]$, the substitution product models functor
composition, directed containers are exactly (small) categories \cite{AhmanUustalu}, and the whole
package reappears inside Spivak's category $\Poly$ of polynomial functors \cite{NiuSpivak}.

The moment one asks for ``a container in a category $\C$ other than $\Set$,'' the theory forks.
At least four constructions claim the name, and they are genuinely different functors:

\begin{itemize}
\item[(1)] \textbf{External families.} An object of $\Fam(\C^\op)$ is a \emph{set} $S$ of shapes
together with position \emph{objects} $P_s\in\C$; the extension $\ext{S,P}=\coprodd_{s\in S}\ihom{P_s}{-}$
uses the self-enrichment of a closed $\C$ and an \emph{external} set-indexed coproduct. This is the
$\Sigma\Pi\mathcal{V}$ construction of Dorta, Jarvis and Niu \cite{DJN}.
\item[(2)] \textbf{Indexed / dependent-polynomial functors.} A diagram $I\leftarrow P\rightarrow S\rightarrow O$
\emph{internal} to a locally cartesian closed base $\C$, with semantics the composite
$\Sigma\,\Pi\,\Delta$ of slice functors --- the polynomial functors of Gambino and Kock \cite{GambinoKock}.
\item[(3)] \textbf{Fibrational containers.} Positions presented as a fibration over the shapes,
with comprehension; this subsumes (1) and (2) as the family and codomain fibrations respectively, and,
applied to the codomain fibration of $\Set$ itself, produces the bifibration \emph{logic of
containers}.
\item[(4)] \textbf{Subcartesian containers.} Walker's locally subcartesian closed bases \cite{Walker},
where an \emph{affine} base-change $\nabla_f\dashv\boxtimes_f$ replaces the top adjoint of the
$\Sigma\dashv\Delta\dashv\Pi$ chain when the base fails to be locally cartesian closed.
\end{itemize}

Over $\C=\Set$ these are one construction viewed four ways. Over a general base they are four
constructions, and the literature offers no systematic answer to the question that then becomes
unavoidable: \emph{for which bases do they agree, and along what axis do they diverge?}

\subsection{The thesis}
We propose a two-property answer. Call a base \emph{good} if it is both \emph{extensive}
(coproducts are disjoint and universal --- equivalently $\C/(A+B)\simeq \C/A\times\C/B$) and
\emph{locally cartesian closed} (every slice is cartesian closed, so pullback has a right adjoint
$\Pi$). The claim is:

\begin{quote}
\emph{Extensivity and local cartesian closure of the base are the two axes along which the four
notions of container diverge. When the base has both, all four agree; and each approach fails to
reach a base precisely when that base fails one of the two properties that the approach silently
assumes.}
\end{quote}

The sharpest instrument for separating the approaches is a base that fails \emph{both} properties.
The category $\Vect$ of vector spaces over a field $k$ is exactly such a base: it is not extensive
(the canonical comparison $A\amalg B \to A\oplus B$ from the coproduct-of-underlying-sets to the
biproduct is not an isomorphism --- ``$\amalg\subsetneq\oplus$''), and it is not locally cartesian
closed (there is no internal dependent product $\Pi$). We therefore run $\Fam(\Vect^\op)$ --- linear
containers --- as the worked example throughout, and it does the discriminating work summarised in
Table~\ref{tab:main}.

\begin{table}[h]
\centering
\renewcommand{\arraystretch}{1.3}
\begin{tabular}{lccccc}
\toprule
Base $\C$ & extensive? & LCC? & (1) $\Fam(\C^\op)$ & (2) $\Sigma\Pi\Delta$ & (4) subcartesian \\
\midrule
$\Set$, any topos & \checkmark & \checkmark & reaches it & forms & (LSCC) \\
$\Vect$ (fd or gen.) & $\times$ ($\amalg\subsetneq\oplus$) & $\times$ (no $\Pi$) & \textbf{reaches it} & \textbf{cannot form $\Pi$} & \textbf{not LSCC} \\
\bottomrule
\end{tabular}
\medskip
\caption{The discriminator. Over a base that is extensive and locally cartesian closed, all
approaches coincide (the classical convergence: containers $\simeq$ directed containers $\simeq$
$\Poly$-comonoids $\simeq$ categories). Over $\Vect$, which has neither property, only the external
approach (1) survives. Approach (3) sits above the table: it presents (1) and (2) as two fibrations
and locates their divergence exactly where the base fails the two properties. This is the honest
reason $\Fam(\Vect^\op)$ ``needs the first approach.''}
\label{tab:main}
\end{table}

\subsection{The quantitative face}
The discriminator is not merely a table of yes/no answers. The obstructions have exact measures,
and supplying them is the technical contribution of the survey. Three theorems about the external
approach (1), each proved in Section~\ref{sec:external} in our own words, turn the two-axis
classification into precise dichotomies. Write $\ext{-}\colon\Fam(\C^\op)\to \C\text{-}[\C,\C]$ for
the container extension into $\C$-enriched endofunctors.

\begin{theorem}[Fullness; T1, Theorem~\ref{thm:T1}]\label{thm:introT1}
$\ext{-}$ is fully faithful if and only if the monoidal unit $I$ is \emph{connected}, i.e.\ the
functor $\C(I,-)\colon\C\to\Set$ preserves small coproducts. This is \emph{not} extensivity of $\C$:
the base $\Set\times\Set$ is extensive yet its unit is disconnected and $\ext{-}$ is not full, while
over $\Vect$ the underlying-set functor $\C(k,-)$ fails to preserve $\oplus$ --- so the unit is
disconnected in this sense --- and $\ext{-}$ is neither full nor faithful.
\end{theorem}

\begin{theorem}[$\otimes$-closedness; T2, Theorem~\ref{thm:T2}]\label{thm:introT2}
The Dirichlet tensor $\otimes$ on $\Fam(\C^\op)$ is closed if and only if a familial-representability
functor $\Phi$ is representable. Over a cartesian base this always holds; over the linear base it
holds \emph{only} on the corner $\Fam_{\mathrm{fin}}(\Vecfd^\op)$ of finite shape-sets and
finite-dimensional positions, failing over both $\Fam(\Vecfd^\op)$ and $\Fam(\Vect^\op)$ by dual
mechanisms.
\end{theorem}

\begin{theorem}[$\lhd$-left-closedness; T4, Theorem~\ref{thm:T4}]\label{thm:introT4}
The substitution product $\lhd$ on $\Fam(\C^\op)$ is left-closed exactly where it \emph{collapses
onto} $\otimes$: over a base whose positions are dualizable (\emph{tiny}), $\lhd=\otimes$, and the
left internal hom is the Dirichlet hom of Theorem~\ref{thm:introT2}. Consequently $\lhd$ is
left-closed on $\Fam_{\mathrm{fin}}(\Vecfd^\op)$ and obstructed over every extensive base. Extensivity
is \emph{opposed} to $\lhd$-left-closedness: the property whose failure costs fullness (T1) and
$\otimes$-closedness (T2) is exactly what buys $\lhd$-closedness.
\end{theorem}

Theorems~\ref{thm:introT1}--\ref{thm:introT4} are the quantitative shadow of the discriminator: the
fullness and closedness obstructions over $\Vect$ are precisely the base failing the extensive/LCC
hypotheses that approaches (2)--(4) would need at the outset. The inversion in
Theorem~\ref{thm:introT4} --- extensivity as villain for T1/T2 and hero-of-the-obstruction for T4
--- is the survey's most surprising structural fact and the reason the classification is genuinely
two-dimensional rather than a single ``goodness'' scale.

\subsection{What is and is not new}
The tensor $\otimes$ over a general base, a composition product $\lhd_{\mathrm{DJN}}$, and the theorem
that $\lhd_{\mathrm{DJN}}$-comonoids are enriched categories (their Theorem~4.3; their Definition~4.2 is
the definition of enriched cofunctor), are due to Dorta, Jarvis and Niu \cite{DJN}; we cite these
as the definitions our theorems are \emph{about}, and re-prove nothing of theirs.
\footnote{\textbf{Scope correction (2026-08-31).} Under the identification
$\Sigma\Pi\mathcal C=\Fam(\mathcal D^\op)$, $\mathcal D=(\Pi\mathcal C)^\op$, our $\otimes$ agrees with
theirs exactly, but their composition product does \emph{not} agree with ours: at a position
$(i,j\colon A_i\to J)$ their direction object is
$\prod_{a\in A_i}\prod_{b\in B_{ja}}(u_{i,a}\cdot v_{ja,b})$ (their Def.~3.5 / Lemma~3.6) --- the
\emph{outer} $u$ multiplied in --- whereas ours is $\prod_{a\in A_i}\prod_{b\in B_{ja}}v_{ja,b}$. The two
coincide exactly at $\mathcal C=1$. A smallest witness: $\mathcal C=2$ with $\cdot=\wedge$, $e=\top$,
$p=\sum_{i\in 1}\prod_{a\in 1}\bot$ and $q=\sum_{j\in 1}\prod_{b\in 1}\top$ give $\bot$ for theirs and
$\top$ for ours. Moreover our defining property $\ext{p\lhd q}\cong\ext p\circ\ext q$ is unavailable on
$\Sigma\Pi\mathcal C$ for $\mathcal C\neq 1$, since their embedding lands in presheaves on
$\Pi\mathcal C$ rather than in endofunctors. We therefore write $\lhd_{\mathrm{DJN}}$ for theirs and
$\lhd$ for ours, and attribute to \cite{DJN} the general-base construction, not a
composition-representing $\lhd$ over a general base.} The $\Set$ fullness
theorem is Abbott--Altenkirch--Ghani \cite{AAG}; our contribution is to name its hidden hypothesis.
The reconstruction of a container from an abstract familially-representable functor is Diers'
\cite{Diers}, and uses extensivity of the codomain, not of the base --- a distinction Section~\ref{sec:external}
turns into the resolution of the ``extensivity'' folklore. Day convolution is \cite{Day}; polynomial
functors over LCC bases are \cite{GambinoKock}; the composition-monoid refinement of the indexed
corner is De Pascalis--Uustalu--Veltr\`i \cite{DPUV}; the subcartesian weakening is Walker
\cite{Walker}. In Section~\ref{sec:fibrational} the facts that $\Fam$ preserves fibrations and that
$\cod$ is a bifibration are folklore \cite{Hermida,Jacobs}, and the fibrewise-op presentation
$\Cont=\int_{\Set}\cod^\op$ is von Glehn's \cite{vonGlehn}. The new material is the two-axis
discriminator, the census of Section~\ref{sec:landscape} together with the assembly of the
individually-known weakenings (Gambino--Kock, Weber, Walker) into one totally ordered tower and the
placement of a given base on it, the three theorems T1/T2/T4 with their named invariants, their
assembly into a single
classification with $\Fam(\Vect^\op)$ as the forcing example, and --- for approach (3) --- the
explicit assembly of $\Cont(\cod)$ as the bifibration of proof-relevant predicates on positions, the
identification of its fibred quantifiers with the $\All/\Exists$ liftings, and the dualisation reading
that makes the container hyperdoctrine the fibrewise opposite of the standard one.

\subsection{Outline}
Section~\ref{sec:prelim} fixes the base and the two tensors. Section~\ref{sec:landscape} maps the
territory: a census of the bases over which the construction has been run, the weakening tower that
orders those between $\Set$ and $\Vect$, and the organising fact that off the good zone a base fails
for exactly one of two reasons. Section~\ref{sec:external} develops the
external approach (1) and proves T1, T2, T4. Section~\ref{sec:indexed} treats the indexed/dependent-
polynomial approach (2) and shows why it cannot form its semantics over $\Vect$.
Section~\ref{sec:fibrational} treats the fibrational approach (3): first as the referee that presents
(1) and (2) as two fibrations, then as a construction in its own right, the bifibration
$\Cont(\cod)$ whose dualised quantifiers are the logic of containers. Section~\ref{sec:subcartesian} treats Walker's approach (4) and records
the verdict that $\Vect$ is not locally subcartesian closed. Section~\ref{sec:synthesis} assembles
the discriminator; Section~\ref{sec:conclusion} names the single sharpest open problem the survey
surfaces.

%==========================================================================
\section{Preliminaries}\label{sec:prelim}

We collect only what the later sections use. Throughout, $(\C,\otimes,I,\ihom{-}{-})$ is a
\emph{closed symmetric monoidal category with small coproducts} $\coprod$. Closedness gives the
enrichment of $\C$ over itself, so that $\C(A,B)=\C(I,\ihom{A}{B})$ naturally, and the adjunction
$-\otimes Q\dashv \ihom{Q}{-}$.

\begin{definition}[Container over $\C$]\label{def:cont}
$\Fam(\C^\op)$ is the category with objects $(S,(P_s)_{s\in S})$, where $S$ is a set of
\emph{shapes} and each $P_s\in\C$ is a \emph{position object}, and morphisms
$(S,P)\to(T,Q)$ the pairs $(f,(\varphi_s))$ with $f\colon S\to T$ a function and
$\varphi_s\colon Q_{f(s)}\to P_s$ in $\C$ (contravariant on positions). Explicitly,
\[
\Fam(\C^\op)\bigl((S,P),(T,Q)\bigr) \;=\; \prodd_{s\in S}\ \coprodd^{\Set}_{t\in T}\ \C(Q_t,P_s),
\]
a choice of one $t=f(s)$ per shape plus a position map. We write $\langle Z\rangle:=(\{*\},Z)$ for
the one-shape generator at $Z\in\C$; every object is the coproduct
$(S,P)=\coprodd_{s\in S}\langle P_s\rangle$, so $\Fam(\C^\op)$ is the free coproduct completion of
$\C^\op$.
\end{definition}

\begin{definition}[Extension]\label{def:ext}
The \emph{extension} of $(S,P)$ is the $\C$-enriched endofunctor
\[
\ext{S,P} \;:=\; \coprodd_{s\in S}\ \ihom{P_s}{-}\ \colon\ \C\longrightarrow\C,
\]
a coproduct of corepresentables. Over $\Set$ (with $I=1$, $\ihom{P}{X}=X^P$) this is the classical
$\ext{S,P}X=\coprodd_s X^{P_s}$; over $\Vect$ (with $I=k$, $\ihom{P}{W}=\Vect(P,W)$, $\coprod=\oplus$)
it is the linear extension $\ext{S,P}W=\bigoplus_s\Vect(P_s,W)$.
\end{definition}

\begin{definition}[The two tensors]\label{def:tensors}
On $\Fam(\C^\op)$ there are two monoidal structures whose extensions we use.
\begin{enumerate}
\item The \emph{Dirichlet} (parallel) tensor $(S,P)\otimes(T,Q):=(S\times T,(P_s\otimes Q_t))$, with
unit $(\{*\},I)$. Its extension is the Day convolution \cite{Day} of the corepresentables:
$\ext{(S,P)\otimes(T,Q)}=\coprodd_{s,t}\ihom{P_s\otimes Q_t}{-}$, using
$\ihom{P}{-}\ast\ihom{Q}{-}=\ihom{P\otimes Q}{-}$.
\item The \emph{substitution} (composition) product $\lhd$, the non-symmetric monoidal structure
whose extension is functor composition: $\ext{(S,P)\lhd(T,Q)}=\ext{S,P}\circ\ext{T,Q}$
\cite{DJN,NiuSpivak}.
\end{enumerate}
\end{definition}

We use four properties of a base, each recalled with the only fact about it we need.

\begin{definition}\label{def:props}
Let $\C$ be as above.
\begin{itemize}
\item $\C$ is \emph{extensive} if finite coproducts are disjoint and universal; equivalently
$\C/(A+B)\simeq\C/A\times\C/B$.
\item $\C$ is \emph{locally cartesian closed} (LCC) if every slice $\C/X$ is cartesian closed;
equivalently every pullback functor $f^*$ has a right adjoint $\Pi_f$ (dependent product).
\item The monoidal unit $I$ is \emph{connected} if $\C(I,-)\colon\C\to\Set$ preserves small
coproducts.
\item An object $Z\in\C$ is \emph{tiny} if $\ihom{Z}{-}\colon\C\to\C$ preserves small coproducts.
In a closed symmetric monoidal category every \emph{dualizable} object is tiny, since then
$\ihom{Z}{-}=Z^*\otimes(-)$ is a left adjoint. Over $\Vect$, tiny $=$ dualizable $=$
finite-dimensional.
\end{itemize}
\end{definition}

\begin{remark}[The linear base is the double-negative example]\label{rem:vecbad}
$\Vect$ fails all the good properties simultaneously. It is not extensive: the comparison
$A\amalg B\to A\oplus B$ from the disjoint union of underlying sets to the biproduct misses every
genuine linear combination $a+b$, so $\coprod^{\Set}\subsetneq\oplus$. It is not LCC: pullback along
a non-mono has no right adjoint, so there is no dependent product $\Pi$ \cite{GambinoKock}. Its unit
$k$ is a connected \emph{object} but $\C(k,-)$ --- the underlying-set functor --- does not preserve
$\oplus$, so $k$ is \emph{not connected} in the sense of Definition~\ref{def:props}. And its tiny
objects are exactly the finite-dimensional ones, so ``tiny'' is a genuine restriction, not
automatic. Each clause powers one obstruction below.
\end{remark}

%==========================================================================
\section{The landscape of bases}\label{sec:landscape}

Table~\ref{tab:main} named the two endpoints of the story: the good base $\Set$, where all four
approaches coincide, and the doubly-bad base $\Vect$, where only the external one survives. Between
them lies a whole territory of bases over which some author has actually run the polynomial/container
construction, each keeping part of the base structure and losing the rest. Before proving anything we
map that territory, because the map is itself an argument for the thesis: every base that has been
used sits at a definite point on the two axes, and which approach reaches it is decided by which axis
it keeps.

\begin{table}[h]
\centering
\renewcommand{\arraystretch}{1.25}
\begin{tabular}{llccl}
\toprule
Base $\C$ & who ran it & ext.? & LCC? & what the construction reaches \\
\midrule
$\Set$, lextensive topos & \cite{AAG,GambinoKock,vonGlehn} & \checkmark & \checkmark & all four approaches coincide \\
finite-limit category & \cite{ShapiroSpivak} & --- & $\times$ & $\lhd$-comonoids $=$ internal categories \\
pullbacks $+$ exp.\ legs & \cite{WeberPoly} & --- & partial & (2) via distributivity pullbacks \\
LSCC (quantale, affine) & \cite{Walker} & $\times$ & $\times$ & (4) affine $\nabla_f\dashv\boxtimes_f$ \\
any monoidal $\Vgen$ & \cite{DJN} & --- & --- & (1) $\otimes,\lhd$; $\lhd$-comon.\ $=$ enr.\ cats \\
any fibration $q$ & \cite{vonGlehn} & --- & --- & (3) $\Cont(q)=\Fam(q^\op)$ \\
$\Set^I$ (indexed) & \cite{DPUV} & \checkmark & \checkmark & (1)/(2) composition monoids \\
$\Prof$ (gen.\ species) & \cite{FGHW} & --- & --- & cartesian closed bicategory \\
$\Vect$, $R\text{-}\mathrm{Mod}$ & this survey & $\times$ & $\times$ & \textbf{(1) only} (T1, T2, T4) \\
$\Rel$;\ $\Mod,\Poly$ & \cite{NiuSpivak} Ch.\,9 & ? & ? & open \\
\bottomrule
\end{tabular}
\medskip
\caption{The bases over which the polynomial/container construction has been run, ordered by
decreasing base strength. The two centre columns are the discriminating axes of the thesis;
``---'' marks a generality-of-base result, for which the axes are not the natural descriptor, and
``partial'' a per-morphism weakening of $\Pi$. The middle block (finite-limit category through LSCC)
is the weakening tower. $\Vect$ and $R\text{-}\mathrm{Mod}$ sit at the doubly-negative corner reached
by the external approach alone.}
\label{tab:landscape}
\end{table}

Reading down the table, the bases that keep pullbacks but lose full local cartesian closure form a
\emph{totally ordered weakening tower}, each rung smoothing a different seam of the $\Sigma\Pi\Delta$
machinery:
\[
\underbrace{\text{LCCC}}_{\text{\cite{GambinoKock}}}\ \supset\
\underbrace{\text{pullbacks}+\text{exp.\ legs}}_{\text{Weber \cite{WeberPoly}}}\ \supset\
\underbrace{\text{subpullbacks (LSCC)}}_{\text{Walker \cite{Walker}}}\ \supset\ \text{spans}.
\]
Local cartesian closure (Gambino--Kock) asks pullback to have a right adjoint globally; Weber's
\emph{distributivity pullbacks} \cite{WeberPoly} weaken this to a per-morphism condition --- a
canonical comparison $\delta$ is invertible iff the relevant middle leg is exponentiable --- so that
composition of polynomials is defined leg-by-leg rather than by a blanket $\Pi$; Walker's subpullbacks
\cite{Walker} weaken it further to an affine base-change $\nabla_f\dashv\boxtimes_f$; and at the bottom
sit bare spans, with no base-change adjoint at all. The rungs are known individually; what has not been
done, and what the survey needs, is to assemble them into one ordered scale and ask where a given
base --- in particular $\Vect$ --- sits on it. The answer, developed in
Section~\ref{sec:subcartesian}, is that $\Vect$ falls off the bottom: not even the affine rung holds.

The tower refines only one of the two axes --- local cartesian closure, which controls the
substitution product $\lhd$ and the internal homs. The other axis, extensivity, is orthogonal to it
and controls whether the external shape-coproduct is seen correctly, hence fullness of $\ext{-}$. This
gives the organising fact of the whole landscape.

\begin{quote}
\emph{Off the good zone, a base fails to support the classical container theory for exactly one of two
independent reasons: loss of extensivity ($\coprod^{\Set}\subsetneq\oplus$, so shapes are no longer
recoverable and $\ext{-}$ is no longer full) or loss of internal $\Pi$ (no dependent product, so
$\lhd$ and closedness collapse). Toposes lose neither; quantales and affine bases lose only the
second, which is why Walker's affine repair reaches them; additive bases such as $\Vect$ and
$R\text{-}\mathrm{Mod}$ lose both, which is why only the external approach reaches them.}
\end{quote}

\noindent This is the two-column heading of Table~\ref{tab:landscape} promoted to the reading of the
whole table, and it is the qualitative form of the thesis that the three theorems of
Section~\ref{sec:external} make quantitative.

\begin{remark}[Where the two proved theorems sit in Weber's work]\label{rem:weberrefile}
The weakening tower places the substitution obstruction and the tensor obstruction in \emph{different}
parts of Weber's work, and it is worth separating them to forestall a tempting conflation. The
hypothesis of the collapse theorem T4 (Theorem~\ref{thm:T4}) --- that the left, composition-side
positions be tiny (exponentiable) --- is exactly Weber's distributivity-$\delta$ condition of
\cite{WeberPoly}: $\delta$ is invertible iff the middle leg is exponentiable, which under the
extension transports to tininess of the left positions and is precisely what makes $\lhd$
well-defined. The hypothesis of the closedness theorem T2 (Theorem~\ref{thm:T2}) --- that a
functor $\Phi$ be \emph{familially representable} --- is instead a comparison from Weber's
\emph{separate} theory of parametric right adjoints and familial functors \cite{WeberFamilial}, where a
functor is familial iff it is fibrewise a coproduct of representables. These conditions are logically
independent: $\delta$ constrains the left (composition) positions, $\Phi$ the right (tensor) positions
and the target, and one can hold while the other fails --- a finite-dimensional-left,
infinite-dimensional-right datum over $\Vect$ makes $\delta$ invertible while $\Phi$ is unrepresentable,
and the reverse datum does the opposite. So T4 belongs to the distributivity-pullback rung of the
tower, whereas T2 does \emph{not} lie on the tower at all: it lives in the parallel familial-functor
theory. This re-filing is a reading of the two Weber papers against the theorems, not a theorem of its
own, and we mark it as such.
\end{remark}

%==========================================================================
\section{Approach (1): external families}\label{sec:external}

This is the approach that reaches $\Vect$, and the one whose obstructions we can measure. We prove
the three theorems that quantify the discriminator. All statements are at the level of the
\emph{set} of enriched natural transformations, which exists for any closed $\C$ with coproducts; no
completeness of $\C$ is needed.

\subsection{The hom formula}
Everything reduces to one enriched computation and one comparison map.

\begin{lemma}[Enriched co-Yoneda]\label{lem:coyoneda}
For any $\C$-endofunctor $G$ and $P\in\C$, $\Nat(\ihom{P}{-},G)\cong \C(I,GP)$, naturally.
\end{lemma}
\begin{proof}
$\ihom{P}{-}=\C(P,-)$ is corepresentable; enriched Yoneda evaluates a transformation at
$\id_P$, and $\C$-linearity of $G$ on hom-objects recovers it everywhere. Taking points
($\C(I,-)$ of the enriched hom) gives the stated set-level form.
\end{proof}

\begin{lemma}[Transformations out of a coproduct]\label{lem:natcoprod}
$\Nat(\coprodd_s F_s, G)\cong\prodd_s\Nat(F_s,G)$.
\end{lemma}
\begin{proof}
Coproducts of endofunctors are pointwise, and a transformation out of a pointwise coproduct is a
tuple of transformations.
\end{proof}

Combining, for the extensions,
\begin{equation}\label{eq:homformula}
\Nat\bigl(\ext{S,P},\ext{T,Q}\bigr)\;\cong\;\prodd_{s\in S}\ \C\Bigl(I,\ \coprodd^{\C}_{t\in T}\ihom{Q_t}{P_s}\Bigr).
\end{equation}
Comparing \eqref{eq:homformula} with the container hom of Definition~\ref{def:cont}, the action of
$\ext{-}$ on morphisms is, in each shape $s$, the following canonical map.

\begin{definition}[The comparison map]\label{def:gamma}
For a family $(X_t)_{t\in T}$ in $\C$,
\[
\gamma_{(X_t)}\ \colon\ \coprodd^{\Set}_{t}\ \C(I,X_t)\ \longrightarrow\ \C\Bigl(I,\ \coprodd^{\C}_{t}X_t\Bigr),
\qquad (t,\,x\colon I\to X_t)\ \longmapsto\ \mathrm{inj}_t\circ x .
\]
This is the comparison testing whether $\C(I,-)$ preserves the coproduct $\coprod_t X_t$. Under
$\C(Q_t,P_s)=\C(I,\ihom{Q_t}{P_s})$, the $s$-component of $\ext{-}$ on morphisms is exactly
$\gamma_{(\ihom{Q_t}{P_s})_t}$.
\end{definition}

\subsection{T1: fullness is unit-connectedness}\label{ss:T1}

\begin{theorem}[Fullness criterion]\label{thm:T1}
Let $\C$ be closed symmetric monoidal with small coproducts.
\begin{enumerate}
\item If the unit $I$ is connected then $\ext{-}\colon\Fam(\C^\op)\to\C\text{-}[\C,\C]$ is fully
faithful.
\item If $\ext{-}$ is full (resp.\ faithful) then for every object $Z$ and set $T$ the map $\gamma$
is surjective (resp.\ injective) for the constant family $(Z)_{t\in T}$; equivalently $\C(I,-)$
preserves (resp.\ reflects) all copowers $T\cdot Z$. A disconnected unit obstructs fullness.
\item If moreover $\C$ has an object $Z_0$ with $\ihom{-}{Z_0}$ essentially surjective onto the
objects appearing in the relevant coproducts (e.g.\ $\C=\Vecfd$, $Z_0=k$, $\ihom{-}{k}=(-)^*$), then
$\ext{-}$ is fully faithful iff $I$ is connected.
\end{enumerate}
\end{theorem}
\begin{proof}
(1) If $\C(I,-)$ preserves coproducts, $\gamma_{(X_t)}$ is a bijection for \emph{every} family, so
by \eqref{eq:homformula} every shape-component of the action is a bijection.

(2) Take source $(\{s\},Z)$ and target $(T,(I)_{t\in T})$, so $\ihom{I}{Z}\cong Z$ and the family is
constant $Z$; the $s$-component of the action is $\gamma_{(Z)_{t\in T}}\colon T\cdot\C(I,Z)\to
\C(I,T\cdot Z)$. Full forces surjectivity for all $T,Z$ (i.e.\ preservation of copowers), faithful
forces injectivity.

(3) With $Z_0$ as stated, $\{\ihom{Q}{Z_0}:Q\in\C\}$ is, up to isomorphism, all objects, so the
families $(\ihom{Q_t}{Z_0})_t$ range over all set-indexed families; fullness plus faithfulness then
give that $\C(I,-)$ preserves all small coproducts, i.e.\ $I$ is connected. With (1) this is the
stated equivalence.
\end{proof}

The point of Theorem~\ref{thm:T1} is that it separates two properties the $\Set$ case fuses.

\begin{corollary}[$\Set$: recovers Abbott--Altenkirch--Ghani]\label{cor:set}
For $\C=\Set$, $I=1$ and $\Set(1,-)=\id$ preserves coproducts, so $\ext{-}$ is fully faithful ---
the representation theorem of \cite{AAG}. The hidden hypothesis is that the terminal object of
$\Set$ is connected.
\end{corollary}

\begin{corollary}[$\Set\times\Set$: extensive but not full]\label{cor:setset}
$\Set\times\Set$ is (l)extensive, yet its unit $(1,1)$ is disconnected:
\[
\C\bigl((1,1),(X,Y)+(X',Y')\bigr)=(X+X')\times(Y+Y')\ \neq\ X{\times}Y+X'{\times}Y',
\]
the comparison missing the cross terms $(x,y')$. By Theorem~\ref{thm:T1}(2) $\ext{-}$ is not full:
with source $(\{s\},(1,1))$ and target $(\{t_1,t_2\},((1,1),(1,1)))$ there are $2$ container
morphisms but $\C((1,1),(2,2))=4$ natural transformations, the two ``crossed'' transformations
unrealised.
\end{corollary}

\begin{corollary}[$\Vect$: the linear collapse]\label{cor:vec}
For $\C=\Vect$, $I=k$, the functor $\Vect(k,-)$ is the underlying-set functor, which does not
preserve $\oplus$: the comparison $\coprod^{\Set}_t X_t\to\bigoplus_t X_t$ is neither surjective (it
misses cross-summand sums) nor injective (it collapses zeros). By Theorem~\ref{thm:T1}(2,3) with
$Z_0=k$, $\ext{-}$ is neither full nor faithful over $\Vect$, and the failure is exactly the
disconnectedness of $k$. This is the $\coprod^{\Set}\subsetneq\oplus$ of the non-extensivity of
$\Vect$, now placed as one instance of the general criterion.
\end{corollary}

\begin{theorem}[The corrected slogan]\label{thm:slogan}
Extensivity of $\C$ is neither necessary nor sufficient for $\ext{-}$ to be full. The controlling
invariant is connectedness of the monoidal unit. ($\Set\times\Set$ is extensive but not full;
$\Set$ is full not because it is extensive but because $1$ is connected.)
\end{theorem}

\begin{remark}[Why the folklore said ``extensivity'']\label{rem:folklore}
Two distinct theorems have been conflated: (a) full faithfulness of the fixed extension $\ext{-}$
(this section: $I$ connected), and (b) Diers' reconstruction \cite{Diers} of $(S,P)$ from an
abstract familially-representable $F\colon\C\to\Set$ as $S=\pi_0(\mathrm{el}\,F)$, which genuinely
uses extensivity of the \emph{codomain} $\Set$. For $\C=\Set$ base and codomain coincide, so (a) and
(b) fuse and ``extensivity'' reads off both. Off $\Set$ they separate: (a) is governed by the unit of
the enrichment base, (b) by extensivity of the value category. Similarly this is no contradiction
with the fully internal polynomial functors of \cite{GambinoKock}: those never form an external
set-coproduct of internal homs, so their representation theorem is untouched. The external-shape
calculus is a $\Set$-enriched phenomenon unless the base's unit is connected.
\end{remark}

\subsection{T2: closedness of the Dirichlet tensor}\label{ss:T2}

We ask when $(-)\otimes(T,Q)$ has a right adjoint inside $\Fam(\C^\op)$.

\begin{theorem}[Closedness $=$ familial representability]\label{thm:T2}
The following are equivalent.
\begin{enumerate}
\item $\otimes$ is closed: $(-)\otimes(T,Q)$ has a right adjoint for every $(T,Q)$.
\item For all families $(Q_t)_{t\in T}$, $(M_r)_{r\in R}$ the functor
$\Phi(Z)=\prodd_{t}\coprodd_{r}\C(M_r,Z\otimes Q_t)\colon\C\to\Set$ is familially representable, i.e.\
$\Phi\cong\coprodd_u\C(N_u,-)$ for a small family $(N_u)$.
\item For all families the atom $\Theta_{A,Q}(Z)=\prodd_{t}\C(A_t,Z\otimes Q_t)$ is familially
representable.
\end{enumerate}
When they hold, $[(T,Q)\Rightarrow(R,M)]=(U,(N_u))$.
\end{theorem}
\begin{proof}
$(1)\Leftrightarrow(2)$: a right adjoint at $(T,Q),(R,M)$ is an object representing
$G:=\Fam(-\otimes(T,Q),(R,M))\colon\Fam(\C^\op)^\op\to\Set$. Since $\otimes$ distributes over the
shape-coproduct and hom out of a coproduct is a product, $G$ sends coproducts to products, and on the
generators $G(\langle Z\rangle)=\prodd_t\coprodd_r\C(M_r,Z\otimes Q_t)=\Phi(Z)$. By the universal
property of the free coproduct completion (Definition~\ref{def:cont}), such a $G$ is representable
iff its restriction to the generators is, i.e.\ iff $\Phi$ is familially representable; and the
representing family is the internal hom.

$(2)\Leftrightarrow(3)$: for sets, $\prodd_t\coprodd_r X_{t,r}=\coprodd_{\rho\colon T\to R}\prodd_t
X_{t,\rho(t)}$ (no hypothesis on $\C$), whence $\Phi\cong\coprodd_\rho\Theta_{M\circ\rho,Q}$ (the
identity \eqref{eq:star} below). A coproduct of familially representable functors is familially
representable, giving $(3)\Rightarrow(2)$; and each corepresentable is connected in $[\C,\Set]$, so
$\Theta_{A,Q}$ appears as the $\rho=\id$ summand of a familially-representable $\Phi$ and is itself
familially representable, giving $(2)\Rightarrow(3)$.
\end{proof}

The set-distributivity used above,
\begin{equation}\label{eq:star}
\Phi(Z)\ \cong\ \coprodd_{\rho\colon T\to R}\ \Theta_{M\circ\rho,\,Q}(Z),\qquad
\Theta_{M\circ\rho,\,Q}(Z)=\prodd_t\C(M_{\rho(t)},Z\otimes Q_t),
\end{equation}
is the engine of both regimes in which $\otimes$ is closed.

\begin{theorem}[Two sufficient regimes]\label{thm:regimes}
\begin{enumerate}
\item \emph{(Cartesian.)} If $\otimes=\times$ (so $\C$ is cartesian closed), $\otimes$ is closed. The
projection split $\C(A,Z\times Q)\cong\C(A,Z)\times\C(A,Q)$ makes each $\Theta_{A,Q}$ a copower of the
corepresentable at $\coprodd_t A_t$, and
\[
[(T,Q)\Rightarrow(R,M)]=\Bigl(\coprodd_{\rho\colon T\to R}\prodd_t\C(M_{\rho(t)},Q_t),\ \bigl(\coprodd_t M_{\rho(t)}\bigr)_\rho\Bigr),
\]
the classical Dirichlet internal hom of $\Poly$ (\cite{NiuSpivak}, Ex.~4.78) over $\Set$.
\item \emph{(Rigid.)} If every $Q_t$ is dualizable and $\C$ has the coproducts $\coprodd_t A_t\otimes
Q_t^*$, then $\otimes$ is closed, with
\[
[(T,Q)\Rightarrow(R,M)]=\Bigl(R^T,\ (N_\rho)_{\rho\colon T\to R}\Bigr),\qquad N_\rho=\coprodd_t M_{\rho(t)}\otimes Q_t^*.
\]
Here $Z\otimes Q_t\cong\ihom{Q_t^*}{Z}$ gives $\C(A_t,Z\otimes Q_t)\cong\C(A_t\otimes Q_t^*,Z)$, so
$\Theta_{A,Q}$ is a single corepresentable at $\coprodd_t A_t\otimes Q_t^*$.
\end{enumerate}
\end{theorem}

Over $\Vect$ neither the projection split nor the diagonal $\Delta$ is available, so the only route
to representability is the rigid one --- and it is also \emph{necessary}.

\begin{lemma}[Single-factor representability over $\Vect$]\label{lem:vecrep}
For $A,Q\in\Vect$ the functor $\Theta(Z)=\Vect(A,Z\otimes Q)$ is familially representable iff $Q$ is
finite-dimensional, in which case it is corepresentable at $A\otimes Q^*$.
\end{lemma}
\begin{proof}
$(\Leftarrow)$ is Theorem~\ref{thm:regimes}(2) with a single $t$. $(\Rightarrow)$: take $A=k$, so
$\Theta(Z)=|Z\otimes Q|$ (underlying set), and $\Theta(0)=1$ forces a familial representation to be a
single corepresentable $\Vect(N,-)$. If $\dim Q=\infty$, write $Q=k^{(I)}$; then $\Theta(Z)$ is the
finitely-supported functions $I\to Z$, which fails to preserve an infinite product $\prodd_J k$: the
element $(e_{i_j})_{j\in J}$ with the $i_j$ pairwise distinct lies in $\prodd_J\Theta(k)$ but has no
finite-support preimage in $\Theta(\prodd_J k)$. So $\Theta$ is not corepresentable, hence not
familially representable.
\end{proof}

\begin{theorem}[The linear closedness dichotomy]\label{thm:vecdicho}
\begin{enumerate}
\item On $\Fam_{\mathrm{fin}}(\Vecfd^\op)$ (finite shape-sets, finite-dimensional positions),
$\otimes$ is closed, with internal hom $(R^T,(\bigoplus_t M_{\rho(t)}\otimes Q_t^*)_\rho)$, which
stays finite-dimensional.
\item On $\Fam(\Vecfd^\op)$ (arbitrary shape-sets), $\otimes$ is not closed: for $T=\mathbb{N}$, all
$Q_t=k$, $R=\{r\}$, $M_r=k$, one computes $\Phi(Z)=|Z|^{\mathbb{N}}=\Vect(k^{(\mathbb{N})},Z)$, whose
representing object $k^{(\mathbb{N})}$ is infinite-dimensional, so no object of $\Vecfd$ represents
it.
\item On $\Fam(\Vect^\op)$ (infinite-dimensional positions), $\otimes$ is not closed: for a single
$Q=k^{(\mathbb{N})}$, $\Phi(Z)=|Z\otimes Q|$ is not familially representable by
Lemma~\ref{lem:vecrep}.
\end{enumerate}
\end{theorem}

\begin{remark}[The named obstruction, and its independence from T1]\label{rem:T2named}
Closedness needs the hom position $N_\rho=\coprodd_t M_{\rho(t)}\otimes Q_t^*$ to (i) \emph{make
sense} --- each $Q_t$ dualizable, forcing finite dimension over a linear base --- and (ii)
\emph{exist} --- the coproduct present in $\C$. Over $\Vect$ these pull apart the moment $T$ is
infinite: (i) pins positions to $\Vecfd$, (ii) needs a colimit escaping $\Vecfd$. The single
load-bearing conjunct is thus \emph{dualizable-and-summable}, satisfiable over $\Vect$ exactly on the
finite/fd corner. This is a different seam from T1: T1 fullness is the failure of $\C(I,-)$ to
preserve $\oplus$ (a property of the unit), whereas T2 is a cocompleteness failure over $\Vecfd$ and
a dualizability failure over full $\Vect$. Cartesian bases dodge the collision because the diagonal
$\Delta$ does the work of duals at no cost --- which is why $\Set\times\Set$, though not full
(Corollary~\ref{cor:setset}), still has $\otimes$ closed. Fullness and $\otimes$-closedness are
independent axes.
\end{remark}

\subsection{T4: left-closedness of the substitution product}\label{ss:T4}

The substitution product $\lhd$ is non-symmetric, so it has two potential internal homs. Its right
internal hom --- the right adjoint to $(T,Q)\lhd(-)$ --- is the $\lhd$-\emph{coclosure}, which
exists over $\Set$ and gives directed containers (\cite{NiuSpivak}, Prop.~6.57). We study the other
one: the \emph{left} internal hom, right adjoint to $(-)\lhd(T,Q)$, the $\lhd$-\emph{closure}. Over
$\Cont=\Fam(\Set^\op)$ it is obstructed.

\begin{proposition}[The $\Set$ obstruction]\label{prop:setobstr}
Over $\Set$, $(-)\lhd q$ has no right adjoint. Computing the extension of the generator,
$\ext{\langle Z\rangle\lhd q}X=\ihom{Z}{\coprodd_t\ihom{Q_t}{X}}$, and the exponential distributes
over the coproduct,
\begin{equation}\label{eq:dist}
\ihom{Z}{\coprodd_t B_t}=\coprodd_{\tau\colon Z\to T}\prodd_{d\in Z}B_{\tau(d)},
\end{equation}
so $\langle Z\rangle\lhd q$ has shape-set $T^Z$. By the universal property of the free coproduct
completion (Definition~\ref{def:cont}), a right adjoint would be a container representing
$(-)\lhd q$ generator by generator; but the $|T|^{|Z|}$-fold shape-set makes its shape-count grow
double-exponentially ($|T_R([n])|\sim n^{2^n}$ in the free case), whereas every container has a
\emph{polynomial} extension. No container carries such data, so the right adjoint does not exist for
$|T|\ge 2$.
\end{proposition}

The obstruction \eqref{eq:dist} is created entirely by $\ihom{Z}{-}$ failing to preserve
$\coprodd_t$ --- the distributive law of an extensive, cartesian closed base. Tininess removes it.

\begin{proposition}[Collapse $\lhd=\otimes$]\label{prop:collapse}
If every position object of $p=(S,P)$ is tiny, then $p\lhd q=p\otimes q=(S\times T,(P_s\otimes Q_t))$
for every $q=(T,Q)$.
\end{proposition}
\begin{proof}
Tininess replaces \eqref{eq:dist} by genuine coproduct-preservation:
\[
\ext{p\lhd q}X=\coprodd_s\ihom{P_s}{\coprodd_t\ihom{Q_t}{X}}
=\coprodd_{s,t}\ihom{P_s}{\ihom{Q_t}{X}}
=\coprodd_{s,t}\ihom{P_s\otimes Q_t}{X}
=\ext{p\otimes q}X,
\]
using tininess of $P_s$ in the middle step and the tensor--hom adjunction at the end.
\end{proof}

Over $\Vecfd$ all objects are dualizable, hence tiny, so Proposition~\ref{prop:collapse} makes $\lhd$
symmetric and identifies its left internal hom with the Dirichlet hom of
Theorem~\ref{thm:vecdicho}.

\begin{theorem}[$\lhd$-left-closedness over the linear base]\label{thm:T4}
With $\lhd$ as above:
\begin{enumerate}
\item $\Fam_{\mathrm{fin}}(\Vecfd^\op)$ is $\lhd$-left-closed: $\lhd=\otimes$ by
Proposition~\ref{prop:collapse}, and $\otimes$ is closed there by Theorem~\ref{thm:vecdicho}(1). The
left internal hom is
\[
[(T,Q)\lhd(R,M)]=\Bigl(R^T,\ \bigl(\coprodd_t M_{\rho(t)}\otimes Q_t^*\bigr)_{\rho\colon T\to R}\Bigr).
\]
\item $\Fam(\Vecfd^\op)$ is not $\lhd$-left-closed (infinite shape-sets break summability, as in
Theorem~\ref{thm:vecdicho}(2)).
\item $\Fam(\Vect^\op)$ is not $\lhd$-left-closed (infinite-dimensional positions are not tiny, so
the collapse itself fails, as in Theorem~\ref{thm:vecdicho}(3)).
\end{enumerate}
\end{theorem}

\begin{theorem}[Extensivity is opposed to $\lhd$-left-closedness]\label{thm:inversion}
The $\Set$ obstruction (Proposition~\ref{prop:setobstr}) and the $\Vecfd$ repair
(Theorem~\ref{thm:T4}) are two valuations of the one fact --- the distributive law
\eqref{eq:dist} --- with opposite sign.
\begin{itemize}
\item Over an extensive, cartesian closed base, $\ihom{Z}{-}$ explodes $\coprodd_t$ into $T^Z$
branches: $\lhd\neq\otimes$, $\lhd$ is genuinely non-symmetric, and its closure is non-polynomial
(absent). Extensivity \emph{is} the obstruction.
\item Over a linear base, $\ihom{Z}{-}$ is additive and preserves $\coprodd_t$: $\lhd$ collapses onto
the closed, symmetric $\otimes$. Non-extensivity \emph{is} the repair.
\end{itemize}
This inverts T1 and T2, where non-extensivity was the obstruction. The reason: T1/T2 measure whether
the external $\Set$-coproduct on shapes is seen correctly by the base (non-extensivity breaks that),
while $\lhd$-closedness needs the base's internal hom on positions \emph{not} to branch against
coproducts --- which is exactly additivity. The same $\coprod^{\Set}\subsetneq\oplus$ seam that costs
fullness buys $\lhd$-closedness.
\end{theorem}

\begin{table}[h]
\centering
\renewcommand{\arraystretch}{1.3}
\begin{tabular}{lccccc}
\toprule
base $\C$ & extensive? & $\lhd$ vs $\otimes$ & $\lhd$ symmetric? & $\otimes$-closed? & $\lhd$-left-closed? \\
\midrule
$\Set$ (toposes) & \checkmark & $\lhd\neq\otimes$ & no & yes & \textbf{no} \\
$\Vecfd$ (fin.\ corner) & $\times$ & $\lhd=\otimes$ & yes & yes & \textbf{yes} \\
\bottomrule
\end{tabular}
\medskip
\caption{The inversion of Theorem~\ref{thm:inversion}: extensivity, the villain of T1/T2, is the
hero of the $\lhd$-closure obstruction.}
\label{tab:inversion}
\end{table}

%==========================================================================
\section{Approach (2): indexed and dependent-polynomial functors}\label{sec:indexed}

The second approach keeps the shapes \emph{internal}. A polynomial in the sense of Gambino and Kock
\cite{GambinoKock} is a diagram
\[
I \xleftarrow{\ s\ } P \xrightarrow{\ f\ } S \xrightarrow{\ t\ } O
\]
in a base $\C$, and its extension is the composite of slice functors
\begin{equation}\label{eq:sigmapidelta}
\C/I \xrightarrow{\ \Delta_s\ } \C/P \xrightarrow{\ \Pi_f\ } \C/S \xrightarrow{\ \Sigma_t\ } \C/O,
\end{equation}
where $\Delta_s=s^*$ is pullback (reindexing), $\Pi_f$ the dependent product (right adjoint to
$f^*$), and $\Sigma_t$ the dependent sum (left adjoint to $t^*$). Setting $I=O=1$ collapses the outer
slices to $\C$ itself and \eqref{eq:sigmapidelta} becomes an endofunctor $\C\to\C$; a ``container in
$\C$'' in this sense is then an internal map $P\to S$, with extension $X\mapsto\Sigma_S\Pi_P\,X$. Over
$\Set$ this is $\coprodd_{s\in S}X^{P_s}$ --- the classical container extension, and the reason the
two approaches coincide there.

\begin{remark}[The load-bearing hypothesis]\label{rem:lcc}
The functor $\Pi_f$ exists precisely when $\C$ is locally cartesian closed. This is not a
convenience: without it the middle arrow of \eqref{eq:sigmapidelta} is undefined, and the approach
cannot even \emph{form} its semantics --- there is no endofunctor to speak of. Over an extensive LCC
base the resulting category of polynomials is itself extensive and reproduces the whole classical
picture, which is Neil Ghani's expectation and is correct for $\Set$, any topos, and any LCC
pretopos.
\end{remark}

\begin{proposition}[The approach cannot reach $\Vect$]\label{prop:novecpoly}
$\Vect$ is not locally cartesian closed (Remark~\ref{rem:vecbad}), so the dependent product $\Pi_f$
in \eqref{eq:sigmapidelta} does not exist. Approach (2) therefore cannot form its semantics over
$\Vect$: there is no $\Sigma\Pi\Delta$ polynomial functor associated to an internal map of vector
spaces. This is a \emph{negative} result --- not that the polynomial functors over $\Vect$ are
badly behaved, but that they are not defined.
\end{proposition}

This is the first half of the discriminator. Where approach (1) sees $\Fam(\Vect^\op)$ directly ---
its external set-coproduct never asks $\C$ for an internal $\Pi$ --- approach (2) is blocked at the
definition. De Pascalis, Uustalu and Veltr\`i \cite{DPUV} give the composition-monoid refinement of
the single-index corner of this approach (a common index, base $\Set$, dropping the cartesianness of
\cite{GambinoKock}); their indexed containers are of the same lineage as (2), still over $\Set$, and
we return to them in Section~\ref{sec:conclusion} as the target of the sharpest open problem.

%==========================================================================
\section{Approach (3): the fibrational referee, and the logic of containers}\label{sec:fibrational}

The third approach begins as a referee. Present the position data as a fibration
$p\colon\E\to\mathbb{B}$ with shapes in the base $\mathbb{B}$ and positions in the fibres, and give
container semantics either by comprehension down to the base or directly in the total category. The
two earlier approaches are then two instances of one scheme:

\begin{itemize}
\item approach (1) is the \emph{family fibration} $\Fam(\C)\to\Set$: shapes index an external set,
fibres are objects of $\C$;
\item approach (2) is the \emph{codomain fibration} $\cod\colon\Arr{\C}\to\C$ of an LCC base:
shapes and positions are both internal, and $\Pi$ is the fibrewise right adjoint that the codomain
fibration of an LCC category carries.
\end{itemize}

\begin{remark}[Where the fibrations diverge]\label{rem:fibdiverge}
The question ``which approach'' becomes ``which fibration models your positions.'' Over an extensive
LCC base the family and codomain fibrations agree, and the choice is presentational. Over a base that
fails one property they genuinely diverge: $\Vect$ has a family fibration $\Fam(\Vect)\to\Set$ (so
approach (1) is available) but no codomain fibration with dependent products (so approach (2)'s
$\Pi$ is absent). The fibrational view thus reads the discriminator as a statement about which
fibration exists.
\end{remark}

But the fibrational approach is more than a referee. Applying the container construction
$\Cont(-):=\Fam((-)^\op)$ to the codomain fibration of $\Set$ itself produces a fibration
\emph{of containers over containers} --- a bifibration whose fibres are proof-relevant predicates on
positions. This is the ``logic of containers,'' and it is where the predicate structure that
approaches (1) and (2) merely presuppose becomes an explicit calculus. We develop it now; it is the
positive content of the third approach and, at the level of the quantifiers, it recovers the
$\All$ and $\Exists$ predicate liftings on positions --- the fibred dependent sum and product.

\subsection{$\Fam$ preserves fibrations}\label{ss:fampreserves}
The construction rests on a single structural lemma. Recall $\Fam(\D)$ has objects $(I,(X_i)_{i\in I})$
with $I\in\Set$, and morphisms $(u,(f_i))\colon(I,(X_i))\to(J,(Y_j))$ a function $u\colon I\to J$ with
$f_i\colon X_i\to Y_{u(i)}$; for $p\colon\E\to\mathbb{B}$ set $\Fam(p)(I,(X_i))=(I,(pX_i))$.

\begin{lemma}[Componentwise cartesianness]\label{lem:fampres}
Let $p\colon\E\to\mathbb{B}$ be any functor. A morphism $(u,(f_i))$ of $\Fam(\E)$ is
$\Fam(p)$-cartesian if and only if each $f_i$ is $p$-cartesian. Consequently, if $p$ is a fibration
(resp.\ opfibration, bifibration) then so is $\Fam(p)$, with (op)reindexing computed componentwise.
\end{lemma}
\begin{proof}[Proof sketch]
Because $\Fam(\E)$ is the free coproduct completion, a morphism carries no cross-index datum beyond
the reindexing function $u$; a factorisation problem for $(u,(f_i))$ therefore decomposes into one
independent factorisation problem per index $i$, each solved uniquely exactly when $f_i$ is
$p$-cartesian. The cartesian lift of a base morphism $(u,(\psi_i))$ into $(J,(Y_j))$ is
$(u,(\chi_i))$, where $\chi_i$ is the $p$-cartesian lift of $\psi_i$ and the position object is
$X_i:=\psi_i^*Y_{u(i)}$. The opcartesian statement is the same argument reversed. The verification is a
routine diagram chase; the fact is folklore \cite[Ch.~1]{Jacobs}, \cite{Hermida}.
\end{proof}

\begin{theorem}[The bifibration of containers]\label{thm:contcod}
The codomain functor $\cod\colon\Arr{\Set}\to\Set$ is a bifibration (opreindexing $\Sigma_\beta$ is
post-composition; reindexing $\beta^*$ is pullback, available because $\Set$ has pullbacks). Hence
$\cod^\op\colon\Arr{\Set}{}^\op\to\Set^\op$ is a bifibration, and by Lemma~\ref{lem:fampres}
\[
\Cont(\cod)\;:=\;\Fam(\cod^\op)\;\colon\;\Cont(\Arr{\Set})\longrightarrow\Cont(\Set)
\]
is a bifibration. Its fibre over a container $(S,(P_s)_{s\in S})$ is
\[
\prodd_{s\in S}\bigl(\Set/P_s\bigr)^{\op},
\]
the \emph{fibrewise opposite} of the slices: an object assigns to each shape $s$ and each position
$b\in P_s$ a witness set, i.e.\ a \emph{proof-relevant predicate on positions}.
\end{theorem}
\begin{proof}[Proof sketch]
$\cod$ is Streicher's fundamental fibration; the op of a bifibration is a bifibration, so
$\cod^\op$ is one and Lemma~\ref{lem:fampres} lifts it through $\Fam$. An object of
$\Cont(\Arr{\Set})$ is $(S,(f_s\colon A_s\to B_s))$ and $\Cont(\cod)$ returns $(S,(B_s))$; a vertical
morphism over $\id_{(S,(P_s))}$ unwinds to a $\Set/P_s$-morphism in the \emph{backward} direction for
each $s$, whence the fibrewise op. This is von Glehn's presentation $\Cont=\int_{\Set}\cod^\op$
\cite{vonGlehn}.
\end{proof}

The fibrewise op is not an accident of bookkeeping: it is forced by the contravariance of container
morphisms on positions, and it is what makes the resulting logic \emph{dual} to the familiar one.

\subsection{The quantifiers are the $\mathsf{All}/\mathsf{Exists}$ liftings}\label{ss:quantifiers}
Fix a container $(S,(P_s))$. Between the trivial predicate $(S,(1)_s)$ and $(S,(P_s))$ there is one
canonical base morphism, the \emph{collapse} $\eta\colon(S,(1))\to(S,(P_s))$ carrying the backward
position map ${!}\colon P_s\to 1$ (the reverse direction would require a choice of point and is not
canonical). Because $\Set$ is locally cartesian closed, collapse along $!\colon P_s\to 1$ gives the
standard adjoint string in each slice,
\[
\Sigma_{!}\ \dashv\ {!}^*\ \dashv\ \Pi_{!}
\qquad(\ \exists\ \dashv\ \text{weakening}\ \dashv\ \forall\ ),
\]
with $\Sigma_{!}(f\colon A\to P_s)=A$ the total space, ${!}^*Z=(P_s\times Z\to P_s)$ the constant
family, and $\Pi_{!}(f)=\prodd_{b}f^{-1}(b)$ the set of sections. Transporting this string through the
fibrewise op and $\Fam$ (an adjunction $L\dashv R$ becomes $R^\op\dashv L^\op$) yields the fibred
quantifiers of $\Cont(\cod)$ along $\eta$.

\begin{theorem}[Quantifiers $=$ liftings]\label{thm:quantAE}
Along the collapse $\eta$, the bifibration $\Cont(\cod)$ carries a two-sided adjoint string
\[
\All\ \dashv\ \Delta_c\ \dashv\ \Exists,
\]
in which the cartesian reindexing $\eta^*=(\Sigma_{!})^\op$ is $\Exists$ (the total space), the
opreindexing $\eta_!=({!}^*)^\op=\Delta_c$ is
weakening, and $\All=(\Pi_{!})^\op$ is the fibrewise product of witnesses. These are the
$\All$ and $\Exists$ predicate liftings on positions: $\Exists$ is the dependent
sum $\Sigma_{b}$ over positions, $\All$ the dependent product $\prodd_{b}$.
\end{theorem}

\begin{remark}[The reindexing is $(\Sigma_\rho)^\op$, not $\rho^*$]\label{rem:trap}
It is tempting to guess that reindexing along a container morphism
$(u,(\rho_s))$ is pullback $\rho_s^*$. It is not: by Lemma~\ref{lem:fampres} the cartesian lift is
computed component\-wise from the $\cod^\op$-cartesian lift, which by op-duality is the
$\cod$-\emph{co}cartesian lift $\Sigma_{\rho_s}$ dualised. So the cartesian reindexing of
$\Cont(\cod)$ is post-composition $(\Sigma_\rho)^\op$, sending a predicate $g\colon C\to P'_{u(s)}$ to
$\rho_s\circ g\colon C\to P_s$. Pullback $\rho_s^*$ reappears, but as the \emph{op}reindexing
(weakening) of the opfibration --- the left adjoint, not the cartesian direction.
\end{remark}

\subsection{The dualisation: containers are the fibrewise-opposite hyperdoctrine}\label{ss:dual}
The adjoint string of Theorem~\ref{thm:quantAE} is the fibrewise opposite of the $\Set$ string
$\Sigma_{!}\dashv{!}^*\dashv\Pi_{!}$. This one fact reorganises the whole predicate logic.

\begin{theorem}[Dualisation]\label{thm:dualisation}
Relative to the common weakening $\Delta_c$, the container logic's existential (left adjoint) is
$\All$, built from $\Set$'s dependent \emph{product} $\Pi$, and its universal (right adjoint) is
$\Exists$, built from $\Set$'s dependent \emph{sum} $\Sigma$: container-$\exists$ is $\Set$-$\forall$
and container-$\forall$ is $\Set$-$\exists$. In each fibre $(\Set/P_s)^\op$ the fibre product
(container $\wedge$) is disjoint union of witness sets, the fibre coproduct (container $\vee$) is the
fibrewise product of witnesses, $\top$ is the empty-witness predicate $\emptyset\to P_s$, and $\bot$
is the full predicate $\id\colon P_s\to P_s$. Hence the hyperdoctrine of containers is the
\emph{fibrewise opposite} of the standard $\Set$-family hyperdoctrine, with $\exists\leftrightarrow
\forall$, $\wedge\leftrightarrow\vee$, $\top\leftrightarrow\bot$ all exchanged, and each fibre is a
\emph{co-topos} whose internal logic is co-intuitionistic (subtractive).
\end{theorem}

The names $\All$ and $\Exists$ track the $\Set$ operation used to build the lifting
($\Pi$ resp.\ $\Sigma$); Theorem~\ref{thm:dualisation} shows their \emph{logical role} in the
container hyperdoctrine is the reverse. The subtractive flavour is genuine only at the proof-relevant
level: the crude Boolean truncation $\Sub(P_s)=2^{P_s}$ is self-dual and hides it, whereas the
witness-level subobject fibration $\Sub(\Arr{\Set})\to\Set$ retains it. Beck--Chevalley and Frobenius
reciprocity hold, dualised: for squares in $\Cont(\Set)$ that are componentwise the op-image of $\Set$
pullback squares, $\Exists$ and $\All$ commute with substitution, and $\All$ satisfies Frobenius with
$\wedge$ replaced by the fibrewise coproduct $\vee$ (co-Frobenius).

\begin{remark}[Scope of the logic]\label{rem:logicscope}
Two honest limitations. First, Theorems~\ref{thm:quantAE}--\ref{thm:dualisation} describe the
\emph{position-level} (fibrewise) logic only; the full first-order logic over $\Cont(\Set)$ also needs
quantification over shapes, which is the $\Fam$ coproduct-completion structure, and we do not verify
here that the combined shape-and-position hyperdoctrine satisfies Beck--Chevalley and Frobenius
jointly. Second, the Beck--Chevalley squares are identified only as the op-image of $\Set$ pullback
squares; an intrinsic characterisation in terms of container pullbacks is left open. Neither gap
affects Theorems~\ref{thm:contcod}--\ref{thm:dualisation} as stated.
\end{remark}

\begin{remark}[Enriched Yoneda and the set-level T1]\label{rem:enriched}
Presenting each object of $\Fam(\C^\op)$ as a presheaf $\C\to\Set$ invites an \emph{enriched}-Yoneda
statement of full faithfulness. Our proof of Theorem~\ref{thm:T1} is deliberately set-level: it
factors through the single comparison map $\gamma$ (Definition~\ref{def:gamma}) and needs no
enrichment of the ambient functor category. The enriched statement --- that $\ext{-}$ is a fully
faithful $\C$-functor --- is the same content repackaged: enriched full faithfulness is exactly the
family of isomorphisms \eqref{eq:homformula}, whose points are the $\gamma$-bijections. The set-level
route is preferable here because the invariant it isolates (connectedness of the unit) is a property
of $\C(I,-)$, not of the enrichment, and this is what makes the $\Set\times\Set$ counterexample
legible.
\end{remark}

For the discriminator, approach (3) plays two roles at once. As referee it reads ``which approach'' as
``which fibration models your positions,'' and the fibrations part company precisely at the
non-extensive, non-LCC bases (Remark~\ref{rem:fibdiverge}). As a construction in its own right over
$\Set$ it supplies the predicate calculus --- the bifibration $\Cont(\cod)$ and its dualised
quantifiers --- that names the logical structure the container of positions carries, and identifies it
with von Glehn's fibrewise op and the $\mathsf{All}/\mathsf{Exists}$ liftings.

%==========================================================================
\section{Approach (4): the subcartesian weakening}\label{sec:subcartesian}

One might hope to rescue approach (2) over $\Vect$ by weakening the requirement that costs it $\Pi$.
Walker's \emph{locally subcartesian closed} (LSCC) categories \cite{Walker} do exactly this, along a
lineage (Weber distributivity pullbacks, Street protocalibrations) independent of Gambino--Kock. The
idea is to equip a category having pullbacks with a coherent choice of \emph{subpullbacks} --- chosen
subobjects of the genuine pullbacks --- from which one builds \emph{affine} base-change functors
$\nabla_f$ whose right adjoints $\boxtimes_f$ (the ``dependent subproduct'') can exist even when the
top adjoint $\Delta_f\dashv\Pi_f$ of the LCC chain does not. Each slice then carries a subcartesian
tensor $f\otimes_X g=\Sigma_g\nabla_g(f)$ that is right-closed without being cartesian closed; the
motivating example is the Lawvere quantale, with the additive, non-cartesian slice tensor $A+B-X$.

This is a genuine fourth weakening, and its additive slice tensor is structurally close to the
$\otimes/\oplus$ shape of $\Vect$. So it sharpens the discriminator: does $\Vect$, though it fails
full LCC, carry a coherent subpullback structure making it LSCC?

\begin{remark}[Verdict: $\Vect$ is not locally subcartesian closed]\label{rem:walkerverdict}
Reading Walker's definitions against $\Vect$ gives a negative answer, on three load-bearing points
\cite[Def.~4.0.1, Rem.~4.0.2, Rem.~4.0.4, fn.~15]{Walker}.
\begin{enumerate}
\item \emph{Pullbacks are mandatory} (Rem.~4.0.4): the construction uses Dubuc's adjoint-triangle
theorem, which needs each $\Sigma_f$ to have a right adjoint, so $\Sigma_f\dashv\Delta_f$ always
holds and only the top adjoint is weakened. The chain does not fully fail; just its top.
\item \emph{Closure demands genuine right adjoints} $\boxtimes_p$ (Def.~4.0.1): the slice tensor must
be \emph{right closed}, i.e.\ each affine base-change must have an actual right adjoint. The
``unique-if-exists'' relaxation applies only to the subpullback (the $\nabla$ side), not to the
closure step, so it gives no back door.
\item \emph{The nearest cousin is filed as obstructed} (Rem.~4.0.2, fn.~15): the category of
\emph{affine spaces} is ``somewhat close to working'' --- it has a canonical $\boxtimes_{X\to 1}$
because slice morphisms are closed under affine combinations --- but general $\boxtimes_p$ fails, the
obstruction being \emph{cross-fiber combinations}. Walker ``mostly ignores such cases as they lack
the desired polynomial structure.''
\end{enumerate}
The cross-fiber obstruction of (3) is exactly the $\coprod^{\Set}\subsetneq\oplus$ biproduct collapse
of $\Vect$ in disguise: linear combinations that cross fibres are precisely what the biproduct adds
and the coproduct-of-underlying-sets lacks. Walker's polynomial biequivalence (Thm.~5.2.8) is a
recognition theorem \emph{assuming} local subcartesian closure, not a criterion certifying it for a
given non-LCC base, so it offers no route around the obstruction.
\end{remark}

\begin{remark}[Status of the verdict]\label{rem:walkerstatus}
Remark~\ref{rem:walkerverdict} is a reading of Walker's preprint checked against the structure of
$\Vect$, not a theorem with a proof of our own; we state it as an understanding result. Its role in
the survey is to close the last escape hatch: $\Vect$ fails not only extensivity and local cartesian
closure but even Walker's weaker affine substitute. The discriminator gap is therefore not ``LCC or
not'' but the deeper ``no affine base-change survives at all.''
\end{remark}

%==========================================================================
\section{Synthesis: the discriminator}\label{sec:synthesis}

We can now state the classification the three (four) approaches obey.

\begin{theorem}[The two-axis discriminator]\label{thm:discriminator}
Let $\C$ be a closed symmetric monoidal base.
\begin{enumerate}
\item If $\C$ is extensive and locally cartesian closed, all four approaches agree: the external
family, indexed-polynomial, family-fibration and codomain-fibration constructions coincide, and one
recovers the classical convergence --- containers $\simeq$ directed containers $\simeq$
$\Poly$-comonoids $\simeq$ (small) categories.
\item If $\C$ fails local cartesian closure, approach (2) cannot form its $\Sigma\Pi\Delta$ semantics
(Proposition~\ref{prop:novecpoly}); if $\C$ additionally fails Walker's subpullback closure, approach
(4) cannot either (Remark~\ref{rem:walkerverdict}).
\item The external approach (1) reaches any closed $\C$ with coproducts, using external set-indexing
and never asking $\C$ for an internal $\Pi$ or a well-behaved $\coprod$. Its obstructions are not to
\emph{existence} but to \emph{good behaviour}, and they are measured exactly:
\begin{itemize}
\item full faithfulness fails iff the unit is disconnected (Theorem~\ref{thm:T1});
\item $\otimes$-closedness fails off the cartesian bases and the finite-fd linear corner
(Theorem~\ref{thm:vecdicho});
\item $\lhd$-left-closedness holds only where $\lhd$ collapses onto $\otimes$, i.e.\ on the same
finite-fd corner, and is obstructed on every extensive base (Theorems~\ref{thm:T4},
\ref{thm:inversion}).
\end{itemize}
\end{enumerate}
The base $\Vect$ realises the corner that has neither property: only approach (1) reaches
$\Fam(\Vect^\op)$, and the three obstructions of (3) are the quantitative face of the two failed
hypotheses.
\end{theorem}

The classification is genuinely two-dimensional. Extensivity and local cartesian closure are
independent (Table~\ref{tab:main} has a base failing each), and they act on the container structures
with \emph{opposite} sign on the two tensors: non-extensivity breaks fullness and
$\otimes$-closedness (T1, T2) but repairs $\lhd$-closedness (T4), because the first two ask the base
to see an external coproduct correctly while the third asks its internal hom not to branch against
coproducts. A single ``goodness'' scale cannot express this; the two axes can. $\Fam(\Vect^\op)$ is
the object that makes the two-dimensionality visible, because it is the corner where both axes are at
their extreme.

%==========================================================================
\section{Conclusion and the sharpest open problem}\label{sec:conclusion}

Four constructions answer to ``a container in a category,'' and they diverge along two independent
properties of the base: extensivity and local cartesian closure. When the base has both, the four
agree and the classical theory is recovered. When the base has neither --- the linear base $\Vect$
--- only the external-family construction survives, and its survival is exactly measured by three
dichotomies: full faithfulness tracks connectedness of the unit, $\otimes$-closedness tracks a
dualizable-and-summable condition met only on the finite-dimensional corner, and $\lhd$-left-closedness
holds precisely where linearity collapses the substitution product onto the Dirichlet tensor. The last
inverts the first two: the non-extensivity that costs fullness and $\otimes$-closedness is what buys
$\lhd$-closedness.

The fibrational approach (3), beyond adjudicating the other two, yields a second thread: over $\Set$
it is the bifibration $\Cont(\cod)=\Fam(\cod^\op)$ of proof-relevant predicates on positions, whose
fibred quantifiers are the $\All$ and $\Exists$ liftings and whose hyperdoctrine is the fibrewise
opposite of the standard one --- container-$\exists$ is $\Set$-$\forall$, and each fibre is a co-topos.
We developed this only at the position level; its shape-level quantifiers, the joint Beck--Chevalley
and Frobenius laws, and the intrinsic class of Beck--Chevalley squares (Remark~\ref{rem:logicscope})
remain open, and are the natural next targets on that thread.

We close by naming the single sharpest problem the survey surfaces, to seed the next proof.

\begin{question}[Does the branching hierarchy lift index-wise?]\label{q:main}
De Pascalis, Uustalu and Veltr\`i \cite{DPUV} classify the monoid structures on \emph{indexed}
containers over $\Set$, dropping the cartesianness of \cite{GambinoKock} and separating a cartesian
from a general regime. Over the single index their composition monad is exactly the substitution
monoid $M\lhd(-)$ studied here. The external approach carries a finer, four-level branching hierarchy
--- $\lambda$-invertible $\subsetneq$ non-branching $\subsetneq$ cartesian $\subsetneq$
$\Pi$-Mendler --- for the liftings of $M\lhd(-)$. \emph{Does this four-level chain lift index-wise to
the indexed containers $\mathrm{IC}_I$ of \cite{DPUV}, refining their cartesian/general dichotomy into
a four-step tower over each index?} A positive answer would join the external approach (1) to the
indexed approach (2) at the level of composition structure, exactly where Section~\ref{sec:indexed}
left them disjoint, and would be the first structural bridge between the two over a base richer than
$\Set$.
\end{question}

A second, more speculative target is whether Walker's bicategory of subcartesian polynomials
(\cite{Walker}, Thm.~5.2.8), built from Street-style spans, is the same object as the external
$\otimes/\lhd$ constructions of Section~\ref{sec:external}, built from family-style external
coproducts, or genuinely disjoint --- the direct overlap check for the container-over-$\Vect$
programme. We record it as secondary; Question~\ref{q:main} is the one to prove first.

%==========================================================================
\section*{Acknowledgements}
This survey was written on a directive of Neil Ghani (Kodamai) to map the approaches to containers
over a base, with $\Fam(\Vect^\op)$ as the running example. The tensor over a general base, and the
weighted composition product $\lhd_{\mathrm{DJN}}$ from which our $\lhd$ is distinguished in
Section~2, are those of Dorta, Jarvis and Niu; the errors of framing are the author's.

%==========================================================================
\begin{thebibliography}{99}

\bibitem{AAG} M.~Abbott, T.~Altenkirch, N.~Ghani.
\emph{Containers: constructing strictly positive types.}
Theoretical Computer Science \textbf{342} (2005), 3--27. (FoSSaCS 2003.)

\bibitem{AhmanUustalu} D.~Ahman, T.~Uustalu.
\emph{Directed containers as categories.}
EPTCS \textbf{207} (2016), 89--98.

\bibitem{Day} B.~Day.
\emph{On closed categories of functors.}
In: Reports of the Midwest Category Seminar IV, Lecture Notes in Mathematics \textbf{137},
Springer, 1970, 1--38.

\bibitem{Diers} Y.~Diers.
\emph{Cat\'egories localisables.}
Th\`ese de doctorat d'\'etat, Universit\'e Paris VI, 1977. (Familial representability.)

\bibitem{DJN} S.~Dorta, A.~Jarvis, N.~Niu.
\emph{Monoidal structures on generalized polynomial categories.}
arXiv:2305.05655.

\bibitem{DPUV} G.~De Pascalis, T.~Uustalu, N.~Veltr\`i.
\emph{Monoid structures on indexed containers.}
arXiv:2509.25879.

\bibitem{FGHW} M.~Fiore, N.~Gambino, M.~Hyland, G.~Winskel.
\emph{The cartesian closed bicategory of generalised species of structures.}
Journal of the London Mathematical Society \textbf{77} (2008), 203--220.

\bibitem{GambinoKock} N.~Gambino, J.~Kock.
\emph{Polynomial functors and polynomial monads.}
Mathematical Proceedings of the Cambridge Philosophical Society \textbf{154} (2013), 153--192.
arXiv:0906.4931.

\bibitem{Hermida} C.~Hermida.
\emph{Fibrations, logical predicates and indeterminates.}
PhD thesis, University of Edinburgh, 1993. (Fam preserves fibrations.)

\bibitem{Jacobs} B.~Jacobs.
\emph{Categorical Logic and Type Theory.}
Studies in Logic and the Foundations of Mathematics \textbf{141}, North-Holland, 1999.

\bibitem{NiuSpivak} N.~Niu, D.~I.~Spivak.
\emph{Polynomial Functors: A Mathematical Theory of Interaction.}
Book draft, 2024.

\bibitem{vonGlehn} T.~von Glehn.
\emph{Polynomials and models of type theory.}
PhD thesis, University of Cambridge, 2015. See also Theory and Applications of Categories
\textbf{33} (2018). (The presentation $\Cont=\int_{\Set}\cod^\op$ with fibrewise op.)

\bibitem{ShapiroSpivak} B.~Shapiro, D.~I.~Spivak.
\emph{Structures on categories of polynomials.}
arXiv:2305.00167.

\bibitem{Walker} C.~Walker.
\emph{Locally subcartesian closed categories.}
arXiv:2607.10242.

\bibitem{WeberPoly} M.~Weber.
\emph{Polynomials in categories with pullbacks.}
Theory and Applications of Categories \textbf{30} (2015), 533--598. arXiv:1106.1983.

\bibitem{WeberFamilial} M.~Weber.
\emph{Familial 2-functors and parametric right adjoints.}
Theory and Applications of Categories \textbf{18} (2007), 665--732.

\end{thebibliography}

\end{document}
